\chapter{Explicit Mean-Field Theory}
\label{ch:mean-field-theory}

This chapter restricts the recognition law of Chapter~\ref{ch:structured-recognition-elbo} while leaving the normalized generative kernel of Chapter~\ref{ch:generative-model} unchanged. The restriction neither replaces an agent by a variational factor nor changes the section-bearing agent ontology. It selects smaller families of probability measures on the same finite-design latent space.

\section{Nested product families on the finite design}

For each agent, collect its coordinates across the declared design as
\begin{equation}
k_i^D=(k_{i,a})_{a=1}^{M},
\qquad
m_i^D=(m_{i,a})_{a=1}^{M}.
\label{eq:mean-field-design-blocks}
\end{equation}
The agent-block family and the finer state--model product family are
\begin{align}
Q_{X,\mathrm{block}}(dY\mid o)
&=\prod_{i=1}^{N}Q_{X,i}(dk_i^D,dm_i^D\mid o),
\label{eq:agent-block-recognition-family}\\
Q_{X,\mathrm{MF}}(dY\mid o)
&=\prod_{i=1}^{N}Q^k_{X,i}(dk_i^D\mid o)Q^m_{X,i}(dm_i^D\mid o).
\label{eq:state-model-recognition-family}
\end{align}
Thus the first family removes dependence between distinct agents but permits dependence between the state and model channels of one agent. The second also removes state--model dependence within each agent. Write \(Q_{X,-(i,k)}\) for the product of every factor in Equation~\eqref{eq:state-model-recognition-family} except \(Q^k_{X,i}\), and write \(Q_{X,-(i,m)}\) for the corresponding complement of \(Q^m_{X,i}\). These finite-design factor measures are distinct from the continuum sections \(q_i\) and \(s_i\). Neither equation requires factorization across the design points inside one displayed factor. A still finer design-point product would be an additional restriction and would require its own declared partition.

For the density-form CAVI equations, assume the finite factor measures are dominated by the corresponding design-product base measures
\begin{equation}
\nu_i^{k,D}:=\bigotimes_{a=1}^{M}\nu_{i,a}^{k},
\qquad
\nu_i^{m,D}:=\bigotimes_{a=1}^{M}\nu_{i,a}^{m},
\label{eq:mean-field-factor-base-measures}
\end{equation}
and define their distinct density symbols by
\begin{align}
Q_{X,i}(dk_i^D,dm_i^D\mid o)
&=\xi^{\mathrm{block}}_{X,i}(k_i^D,m_i^D\mid o)
\left(\nu_i^{k,D}\otimes\nu_i^{m,D}\right)(dk_i^D,dm_i^D),
\label{eq:agent-block-factor-density}\\
Q^k_{X,i}(dk_i^D\mid o)
&=\xi^k_{X,i}(k_i^D\mid o)\nu_i^{k,D}(dk_i^D),
\label{eq:state-factor-density}\\
Q^m_{X,i}(dm_i^D\mid o)
&=\xi^m_{X,i}(m_i^D\mid o)\nu_i^{m,D}(dm_i^D).
\label{eq:model-factor-density}
\end{align}
Thus \(Q_{X,i},Q^k_{X,i},Q^m_{X,i}\) are probability measures, whereas \(\xi^{\mathrm{block}}_{X,i},\xi^k_{X,i},\xi^m_{X,i}\) are their Radon--Nikodym densities. The three explicit Markov-blanket coordinate logarithms below apply to these densities, never to the capital factor measures.

Let \(\mathfrak B\) denote any finite measurable partition of the coordinates of \(Y_D\), and write \(Y=(Y_b,Y_{-b})\) for one block and its complement. A product-family member has density
\begin{equation}
q(y)=q_b(y_b)q_{-b}(y_{-b}),
\qquad
q_{-b}(y_{-b})=\prod_{c\in\mathfrak B\setminus\{b\}}q_c(y_c),
\label{eq:generic-product-recognition-density}
\end{equation}
relative to the corresponding finite product base measure. The product factors are recognition objects only; every expectation below is evaluated against the same fixed \(p_\theta(o,y\mid X)\).

\section{Normalized functional coordinate equation}

Fix every factor except \(q_b\), and define the expected log joint
\begin{equation}
g_b(y_b)
=\E_{Q_{-b}}\left[\log p_\theta(o,y_b,Y_{-b}\mid X)\right].
\label{eq:cavi-expected-log-joint}
\end{equation}
Assume that \(g_b\) has a measurable finite version \(\nu_b\)-almost everywhere, that the product ELBO is well-defined under the componentwise absolute-log-integrability condition used in Equation~\eqref{eq:state-free-energy-factorization}, and that
\begin{equation}
0<Z_b
:=\int \exp\bigl(g_b(u_b)\bigr)\nu_b(du_b)<\infty.
\label{eq:cavi-coordinate-existence}
\end{equation}
These conditions permit Fubini's theorem and define a normalized candidate. If the integral is zero or infinite, or if the expected log joint has no measurable finite version \(\nu_b\)-almost everywhere, the coordinate displayed below need not exist.

For a bounded signed perturbation \(\delta q_b\) with \(\int\delta q_b d\nu_b=0\), the terms of the state ELBO that depend on \(q_b\) are
\begin{equation}
\mathcal J_b(q_b)
=\int q_b(y_b)g_b(y_b)\nu_b(dy_b)
-\int q_b(y_b)\log q_b(y_b)\nu_b(dy_b).
\label{eq:cavi-coordinate-functional}
\end{equation}
Introducing a multiplier \(\lambda_b\) for \(\int q_b d\nu_b=1\), the first variation is
\begin{equation}
\delta\left\{\mathcal J_b(q_b)
+\lambda_b\left(\int q_b d\nu_b-1\right)\right\}
=\int\delta q_b(y_b)
\left[g_b(y_b)-\log q_b(y_b)-1+\lambda_b\right]\nu_b(dy_b).
\label{eq:cavi-constrained-first-variation}
\end{equation}
Stationarity for every admissible perturbation, followed by normalization, gives
\begin{align}
\log q_b^\star(y_b)
&=\E_{Q_{-b}}\left[\log p_\theta(o,Y\mid X)\right]-\log Z_b,
\label{eq:cavi-normalized-coordinate-log}\\
q_b^\star(y_b)
&=\frac{\exp\left(\E_{Q_{-b}}[\log p_\theta(o,Y\mid X)]\right)}
{\displaystyle\int
\exp\left(\E_{Q_{-b}}[\log p_\theta(o,u_b,Y_{-b}\mid X)]\right)
\nu_b(du_b)}.
\label{eq:cavi-normalized-coordinate-density}
\end{align}
The additive constant in Equation~\eqref{eq:cavi-normalized-coordinate-log} is explicitly the negative log integral \(-\log Z_b\), not an unspecified parameter. Equivalently, subtracting the value at \(q_b^\star\) from Equation~\eqref{eq:cavi-coordinate-functional} gives \(-\KL(q_b\Vert q_b^\star)\), so \(q_b^\star\nu_b\) is the unique coordinate-maximizing probability measure whenever the displayed relative entropy is defined; its density version is unique only up to \(\nu_b\)-null sets \citep{Bishop2006,Blei2017,Wainwright2008}.

\section{Complete directed Markov blankets}

Let \(\operatorname{ch}(i)=\{\ell\in V:i\in\operatorname{pa}(\ell)\}\). For each design point, define the receiving log factors using the root conventions of Chapter~\ref{ch:generative-model}:
\begin{align}
\ell_{i,a}^{m}
&=
\begin{cases}
\log p_{\theta,i,a}^{m}(m_{i,a}\mid X),&i\in\mathcal R,\\
\log p_{\theta,i,a}^{m}(m_{i,a}\mid m_{\operatorname{pa}(i),a},X),&i\notin\mathcal R,
\end{cases},
\label{eq:model-receiving-log-factor}\\
\ell_{i,a}^{k}
&=
\begin{cases}
\log p_{\theta,i,a}^{k}(k_{i,a}\mid m_{i,a},X),&i\in\mathcal R,\\
\log p_{\theta,i,a}^{k}(k_{i,a}\mid k_{\operatorname{pa}(i),a},m_{i,a},X),&i\notin\mathcal R,
\end{cases},
\label{eq:state-receiving-log-factor}\\
\ell_{i,a}^{o}
&=\log p_{\theta,i,a}^{o}(o_{i,a}\mid k_{i,a},m_{i,a},X).
\label{eq:observation-log-factor}
\end{align}

In the state--model family of Equation~\eqref{eq:state-model-recognition-family}, the exact coordinate for the complete state factor of agent \(i\) is
\begin{align}
\log \bigl(\xi^k_{X,i}\bigr)^\star(k_i^D\mid o)
={}&\sum_{a=1}^{M}\E_{Q_{X,-(i,k)}}\left[
\ell_{i,a}^{k}+\ell_{i,a}^{o}
+\sum_{\ell\in\operatorname{ch}(i)}
\log p_{\theta,\ell,a}^{k}
(k_{\ell,a}\mid k_{\operatorname{pa}(\ell),a},m_{\ell,a},X)
\right]
-\log Z_i^k.
\label{eq:complete-state-cavi-coordinate}
\end{align}
The child sum is present because \(k_i\) occurs as a conditioning variable in every outgoing state transition. Likewise, the complete model-factor coordinate is
\begin{align}
\log \bigl(\xi^m_{X,i}\bigr)^\star(m_i^D\mid o)
={}&\sum_{a=1}^{M}\E_{Q_{X,-(i,m)}}\left[
\ell_{i,a}^{m}+\ell_{i,a}^{k}+\ell_{i,a}^{o}
+\sum_{\ell\in\operatorname{ch}(i)}
\ell_{\ell,a}^{m}
\right]
-\log Z_i^m.
\label{eq:complete-model-cavi-coordinate}
\end{align}
Here \(\ell_{i,a}^{k}\) is the local state bridge or receiving transition, which depends on \(m_i\), and \(\ell_{i,a}^{o}\) is included because the declared observation kernel depends on both local channels. Every outgoing model transition appears because its parent tuple contains \(m_i\). No outgoing child-state factor in the fixed joint depends on a parent model coordinate; if a different joint introduced such a dependence, that factor would also belong in Equation~\eqref{eq:complete-model-cavi-coordinate}. The constants \(-\log Z_i^k\) and \(-\log Z_i^m\) normalize the complete design-block factors under the existence conditions of Equation~\eqref{eq:cavi-coordinate-existence}.

For the coarser agent block, the corresponding complete coordinate combines both blankets:
\begin{align}
&\log \bigl(\xi^{\mathrm{block}}_{X,i}\bigr)^\star
\left(k_i^D,m_i^D\mid o\right)\notag\\
&\quad={}
\sum_{a=1}^{M}\E_{Q_{X,-i}}\left[
\ell_{i,a}^{m}+\ell_{i,a}^{k}+\ell_{i,a}^{o}
\right.\notag\\
&\left.\qquad+
\sum_{\ell\in\operatorname{ch}(i)}
\left\{
\ell_{\ell,a}^{m}+\ell_{\ell,a}^{k}
\right\}
\right]
-\log Z_i^{\mathrm{block}}.
\label{eq:complete-agent-block-cavi-coordinate}
\end{align}

An update that retains only the receiving factor and perhaps the local observation while omitting either child sum is a receiver-only local surrogate. It is not coordinate-ascent variational inference for the ELBO in Equation~\eqref{eq:state-elbo-definition}. This distinction holds even if the surrogate is numerically stable or shares a fixed point with CAVI in a special parameter regime.

\section{Gaussian reverse-KL optimum and evidence gap}

Assume the exact posterior for the stacked \(d\)-dimensional finite-design latent vector is
\begin{equation}
P_\theta(dy\mid o,X)=\mathcal N(m,J^{-1})(dy),
\qquad
J=J^\top\succ0,
\label{eq:gaussian-exact-posterior}
\end{equation}
and declare a block partition \(\mathfrak B\) with block dimensions \(d_b\) satisfying \(\sum_{b\in\mathfrak B}d_b=d\). Restrict recognition to
\begin{equation}
Q=\mathcal N(\mu,C),
\qquad
C=\operatorname{blockdiag}_{b\in\mathfrak B}(C_b),
\qquad
C_b\in\mathbb S_{++}^{d_b}.
\label{eq:gaussian-block-product-family}
\end{equation}
Direct evaluation of the Gaussian density ratio gives
\begin{equation}
\KL(Q\Vert P_\theta)
=\frac12\left[
\operatorname{tr}(JC)
+(\mu-m)^\top J(\mu-m)
-d-\log\det J-\log\det C
\right].
\label{eq:gaussian-reverse-kl-block}
\end{equation}
The first-principles calculation and matrix derivatives are recorded in Appendix~\ref{sec:gaussian-block-reverse-kl-proof}. Differentiation on the open SPD blocks yields
\begin{equation}
\nabla_\mu\KL=J(\mu-m),
\qquad
\nabla_{C_b}\KL=\frac12\left(J_{bb}-C_b^{-1}\right).
\label{eq:gaussian-reverse-kl-derivatives}
\end{equation}
Strict positive definiteness and strict convexity in each covariance block therefore give the unique optimum
\begin{equation}
\mu^\star=m,
\qquad
C_b^\star=J_{bb}^{-1}.
\label{eq:gaussian-block-reverse-kl-optimum}
\end{equation}
The mean is exact even though the recognition covariance is restricted. The covariance result concerns the diagonal blocks of posterior precision. In general,
\begin{equation}
J_{bb}^{-1}\neq(J^{-1})_{bb},
\label{eq:precision-block-not-marginal-covariance}
\end{equation}
because \((J^{-1})_{bb}\) is the posterior marginal covariance and includes the Schur-complement effect of every other block.

The same conclusion appears from Gaussian CAVI. Writing the posterior information vector as \(h=Jm\), one coordinate has covariance \(J_{bb}^{-1}\) and mean update
\begin{equation}
\mu_b^{\mathrm{new}}
=J_{bb}^{-1}
\left(h_b-\sum_{c\in\mathfrak B\setminus\{b\}}J_{bc}\mu_c\right).
\label{eq:gaussian-cavi-mean-update}
\end{equation}
Every fixed point satisfies \(J\mu=h\), hence \(\mu=m\).

Write \(\mathcal L_{\mathfrak B}^\star=\Lstate(Q_{\mathfrak B}^\star;X)\) for the state ELBO evaluated at this block-Gaussian optimum. At the optimum, \(\operatorname{tr}(JC^\star)=d\) and \(\log\det C^\star=-\sum_b\log\det J_{bb}\). Combining Equation~\eqref{eq:state-evidence-identity} with Equation~\eqref{eq:gaussian-reverse-kl-block} gives
\begin{equation}
\log p_\theta(o\mid X)-\mathcal L_{\mathfrak B}^\star
=\frac12\log
\frac{\displaystyle\prod_{b\in\mathfrak B}\det J_{bb}}
{\det J}.
\label{eq:gaussian-mean-field-evidence-gap}
\end{equation}
The block Fischer determinant inequality for \(J\succ0\) states
\begin{equation}
\det J\leq\prod_{b\in\mathfrak B}\det J_{bb},
\label{eq:block-fischer-determinant-inequality}
\end{equation}
with equality if and only if every off-diagonal block under the declared partition vanishes \citep{HornJohnson2013}. The gap is therefore nonnegative and is zero exactly when \(J\) is block diagonal under \(\mathfrak B\). Appendix~\ref{sec:gaussian-block-reverse-kl-proof} proves the equality condition by successive Schur complements.

For Equation~\eqref{eq:agent-block-recognition-family}, the blocks are \((k_i^D,m_i^D)\), so the criterion removes only cross-agent posterior-precision blocks. For Equation~\eqref{eq:state-model-recognition-family}, the blocks are the separate \(k_i^D\) and \(m_i^D\) coordinates, so the criterion also removes every within-agent state--model precision block. The two determinant gaps refer to different partitions and need not agree.

\begin{figure}[tbp]
\centering
\begin{tikzpicture}[
  >={Latex[length=2.2mm]},
  every node/.style={font=\small},
  matrixcell/.style={draw=darkgray,line width=0.55pt,minimum width=1.16cm,minimum height=0.82cm,anchor=center},
  matrixlabel/.style={font=\bfseries\small,text=primaryblue,align=center},
  restriction/.style={->,line width=0.95pt,draw=criticalred},
  callout/.style={draw=darkgray,line width=0.8pt,rounded corners=2pt,fill=lightgray,align=center,inner xsep=6pt,inner ysep=5pt}
]
\matrix (fullprecision) [
  matrix of math nodes,
  nodes={matrixcell},
  row sep=-\pgflinewidth,
  column sep=-\pgflinewidth,
  left delimiter={[},
  right delimiter={]}
] {
  |[fill=lightblue]| J_{11} & |[fill=lightorange]| J_{12} & |[fill=lightorange]| J_{13} \\
  |[fill=lightorange]| J_{12}^{\mathsf T} & |[fill=lightgreen]| J_{22} & |[fill=lightorange]| J_{23} \\
  |[fill=lightorange]| J_{13}^{\mathsf T} & |[fill=lightorange]| J_{23}^{\mathsf T} & |[fill=lightpurple]| J_{33} \\
};
\node[matrixlabel,above=0.34cm of fullprecision] (fulllabel)
  {exact posterior precision\\\(J=J^{\mathsf T}\succ0\)};

\matrix (restrictedcovariance) [
  matrix of math nodes,
  nodes={matrixcell},
  row sep=-\pgflinewidth,
  column sep=-\pgflinewidth,
  left delimiter={[},
  right delimiter={]},
  right=4.45cm of fullprecision
] {
  |[fill=lightblue]| J_{11}^{-1} & 0 & 0 \\
  0 & |[fill=lightgreen]| J_{22}^{-1} & 0 \\
  0 & 0 & |[fill=lightpurple]| J_{33}^{-1} \\
};
\node[matrixlabel,above=0.34cm of restrictedcovariance] (restrictedlabel)
  {block-product variational covariance\\\(C^\star=\operatorname{blockdiag}_b(J_{bb}^{-1})\)};

\draw[restriction] (fullprecision.east) --
  node[midway,above=0.16cm,font=\scriptsize,align=center,text=criticalred]
  {recognition-family\\restriction}
  node[midway,below=0.16cm,font=\scriptsize,align=center,text=darkgray]
  {removes conditional\\dependencies \(J_{bc}\), \(b\ne c\)}
  (restrictedcovariance.west);

\coordinate (matrixmid) at ($(fullprecision.south)!0.5!(restrictedcovariance.south)$);
\node[callout,below=1.35cm of matrixmid] (gap)
  {\(\displaystyle
    \log p_\theta(o\mid X)-\mathcal L_{\mathfrak B}^{\star}
    =\frac12\log\frac{\prod_{b\in\mathfrak B}\det J_{bb}}{\det J}\geq0
  \)};
\node[below=0.18cm of gap,font=\scriptsize,text=darkgray,align=center]
  {The restriction does not set posterior marginal covariances \((J^{-1})_{bc}\) to zero.};
\end{tikzpicture}

\caption{\textbf{Gaussian mean field restricts conditional dependence, not posterior covariance.} The exact nondegenerate Gaussian posterior has full precision \(J\). Restricting recognition to the declared block-product family gives the exact reverse-KL optimum \(C_b^\star=J_{bb}^{-1}\) within that family and the displayed nonnegative determinant gap. This restriction removes off-diagonal conditional-dependence blocks from the variational precision description; it does not assert that the exact marginal covariance blocks \((J^{-1})_{bc}\) vanish. A block may collect one agent or split one agent's channels, so matrix blocks must not be identified with the agent ontology without the stated partition.}
\label{fig:magent-precision-mean-field}
\end{figure}

\section{Partition-specific zero-coupling controls}

The Gaussian reference joint supplies a direct precision audit. Write
\begin{equation}
T_i^k=A_i^k\Omega_{ij}^k,
\qquad
T_i^m=A_i^m\Omega_{ij}^m,
\qquad
\Lambda_i^k=(Q_i^k)^{-1},
\qquad
\Lambda_i^m=(Q_i^m)^{-1},
\label{eq:gaussian-transition-precision-notation}
\end{equation}
for a nonroot edge \(j\to i\), with the root state precision denoted \(\Lambda_{r,0}^k=(Q_{r,0}^k)^{-1}\). Expanding the Gaussian residuals shows that a nonroot model transition contributes the cross-agent precision block \(-\Lambda_i^mT_i^m\) between \(m_i\) and \(m_j\). A nonroot state transition contributes \(-\Lambda_i^kT_i^k\) between \(k_i\) and \(k_j\), \(-\Lambda_i^kB_i\) between \(k_i\) and \(m_i\), and \((T_i^k)^\top\Lambda_i^kB_i\) between \(k_j\) and \(m_i\), together with their transposes. The local observation contributes
\begin{equation}
H_i^\top(R_i^o)^{-1}L_i,
\label{eq:observation-state-model-cross-precision}
\end{equation}
to the \((k_i,m_i)\) precision block. Appendix~\ref{sec:gaussian-zero-coupling-proof} derives these signs from the quadratic residuals.

For the agent-block control, require the total assembled posterior-precision block between any two distinct agent blocks to vanish at every design point. A transparent sufficient construction sets \(T_i^k=0\) and \(T_i^m=0\) on every directed edge, thereby removing all parent dependence while preserving normalized receiving kernels. More general cancellations are admissible only when the sum of all cross-agent precision contributions is exactly zero. The bridge matrices \(B_i\) and the local observation cross curvature in Equation~\eqref{eq:observation-state-model-cross-precision} may remain. Since every remaining latent quadratic term is contained within one \((k_i^D,m_i^D)\) block, the posterior precision has the form
\begin{equation}
J=\operatorname{blockdiag}_{i=1}^{N}J_{(k_i^D,m_i^D)},
\label{eq:agent-block-zero-coupling-precision}
\end{equation}
the exact posterior belongs to the family in Equation~\eqref{eq:agent-block-recognition-family}, and its determinant gap is zero.

For the fully factorized state--model control, impose the agent-block condition and also require the total \((k_i^D,m_i^D)\) cross-precision block to vanish for every agent. In the displayed Gaussian reference, the local condition at a nonroot is
\begin{equation}
-\Lambda_i^kB_i
+H_i^\top(R_i^o)^{-1}L_i
+J_{k_i m_i}^{\mathrm{other}}=0,
\label{eq:fully-factorized-local-cross-precision-control}
\end{equation}
with \(-\Lambda_{r,0}^kB_r\) replacing the first term at a root. The term \(J_{k_i m_i}^{\mathrm{other}}\) denotes the sum of any additional fixed factors that couple the two local channels; it is zero for the joint in Chapter~\ref{ch:generative-model}. Under these conditions,
\begin{equation}
J=\operatorname{blockdiag}
\left(J_{k_1^D},J_{m_1^D},\ldots,J_{k_N^D},J_{m_N^D}\right),
\label{eq:fully-factorized-zero-coupling-precision}
\end{equation}
so the exact posterior belongs to Equation~\eqref{eq:state-model-recognition-family} and the finer determinant gap is zero. Setting one named edge coefficient to zero does not prove either block-diagonality statement, and a small numerical off-diagonal block does not establish exact equality. The theorem requires the complete assembled precision, including all transition, bridge, observation, and other coupling contributions.

\section{A fixed-source categorical row model}

Let \((\mathsf Y_i,\mathscr Y_i,\nu_i)\) be a state sample space. Fix normalized densities \(f_{ij}\) on this space and a normalized prior row \(\pi_i=(\pi_{ij})_{j\in\mathcal J_i}\) with \(\pi_{ij}>0\). Let \(\varrho_i\) be a normalized density satisfying \(\varrho_i\ll f_{ij}\) and \(\KL(\varrho_i\Vert f_{ij})<\infty\) for every candidate source. Define two normalized joint laws on \(\mathsf Y_i\times\mathcal J_i\) by
\begin{equation}
P_i(y,j)=\pi_{ij}f_{ij}(y),
\qquad
Q_i(y,j)=\beta_{ij}\varrho_i(y),
\qquad
\sum_j\beta_{ij}=1.
\label{eq:fixed-source-row-joints}
\end{equation}
Direct expansion of the joint density ratio gives
\begin{align}
\KL(Q_i\Vert P_i)
&=\sum_j\beta_{ij}
\int \varrho_i(y)
\log\frac{\beta_{ij}\varrho_i(y)}{\pi_{ij}f_{ij}(y)}\nu_i(dy)
\label{eq:fixed-source-row-kl-expansion}\\
&=\sum_j\beta_{ij}\KL(\varrho_i\Vert f_{ij})
+\KL(\beta_i\Vert\pi_i).
\label{eq:fixed-source-row-kl-decomposition}
\end{align}
Minimizing this exact joint KL over the probability simplex and normalizing the Lagrange stationary equation yields
\begin{equation}
\beta_{ij}^\star
=\frac{\pi_{ij}\exp[-\KL(\varrho_i\Vert f_{ij})]}
{\displaystyle\sum_{\ell\in\mathcal J_i}
\pi_{i\ell}\exp[-\KL(\varrho_i\Vert f_{i\ell})]}.
\label{eq:fixed-source-row-optimum}
\end{equation}
This formula is exact for fixed normalized component densities and unit temperature. A nonunit temperature is not obtained by inserting a scale into Equation~\eqref{eq:fixed-source-row-optimum} while claiming the same joint KL. It requires a separately normalized tempered construction, including each component's temperature-dependent normalizer. Refreshing \(f_{ij}\) from live recognition marginals also changes \(P_i\) and therefore changes the model between updates.

The categorical entropy of \(\beta_i\) is an auxiliary source-row entropy on \(\mathcal J_i\). It is distinct from the state entropy of \(\varrho_i\) on \(\mathsf Y_i\), and both are distinct from the configuration entropy of a recognition law on \((\mathsf X,\mathscr X)\). Equation~\eqref{eq:fixed-source-row-kl-decomposition} supplies no license to merge or double count those sectors.
The row-local construction is normalized for one receiver. Multiplying two or more such rows does not by itself define a normalized population joint when their state variables are shared; that would require a separately specified global construction.

\section{A tempered source row and its normalizer}
\label{sec:tempered-source-row}

Equation~\eqref{eq:fixed-source-row-optimum} is exact at unit ordinary-KL scale, so a nonunit temperature is not obtained by inserting a scale into it. This section supplies the separately normalized tempered construction that such a scale requires and records exactly what its normalizer contributes to the row. The device is the power-likelihood construction, whose normalized form is standard \citep{BissiriHolmesWalker2016}.

Retain the setting of Equation~\eqref{eq:fixed-source-row-joints}. For \(\tau>0\) assume
\begin{equation}
0<Z_{ij}(\tau)
:=\int_{\mathsf Y_i}f_{ij}(y)^{1/\tau}\nu_i(dy)<\infty
\qquad\text{for every }j\in\mathcal J_i,
\label{eq:tempered-component-normalizer}
\end{equation}
and define the tempered component and the tempered joint
\begin{equation}
f_{ij}^{(\tau)}:=\frac{f_{ij}^{1/\tau}}{Z_{ij}(\tau)},
\qquad
P_i^{(\tau)}(y,j)=\pi_{ij}f_{ij}^{(\tau)}(y).
\label{eq:tempered-source-joint}
\end{equation}
Each \(f_{ij}^{(\tau)}\) is a normalized density by construction, so \(P_i^{(\tau)}\) is a normalized joint on \(\mathsf Y_i\times\mathcal J_i\) and the decomposition of Equation~\eqref{eq:fixed-source-row-kl-decomposition} applies verbatim with \(f_{ij}\) replaced by \(f_{ij}^{(\tau)}\).

\paragraph{Proposition (tempered row decomposition).}
\label{prop:tempered-row-decomposition}
Assume \(\varrho_i\ll f_{ij}\), that \(\KL(\varrho_i\Vert f_{ij})\) is finite, that the differential entropy \(H(\varrho_i)=-\int\varrho_i\log\varrho_i d\nu_i\) is finite, and that Equation~\eqref{eq:tempered-component-normalizer} holds. Then
\begin{equation}
\KL\left(\varrho_i\middle\Vert f_{ij}^{(\tau)}\right)
=\frac{1}{\tau}\KL(\varrho_i\Vert f_{ij})
+\left(\frac1\tau-1\right)H(\varrho_i)
+\log Z_{ij}(\tau),
\label{eq:tempered-kl-decomposition}
\end{equation}
and consequently the exact simplex minimizer of \(\KL(Q_i\Vert P_i^{(\tau)})\) is
\begin{equation}
\beta_{ij}^{\star(\tau)}
=\frac{\pi_{ij}Z_{ij}(\tau)^{-1}\exp\left[-\KL(\varrho_i\Vert f_{ij})/\tau\right]}
{\displaystyle\sum_{\ell\in\mathcal J_i}
\pi_{i\ell}Z_{i\ell}(\tau)^{-1}\exp\left[-\KL(\varrho_i\Vert f_{i\ell})/\tau\right]}.
\label{eq:tempered-source-row-optimum}
\end{equation}
The proof is in Appendix~\ref{sec:tempered-and-labeled-source-proofs}.

The middle term of Equation~\eqref{eq:tempered-kl-decomposition} carries no source index and therefore cancels under row normalization. The normalizer does not cancel, and its source dependence is the entire difference between Equation~\eqref{eq:tempered-source-row-optimum} and a tempered reading of Equation~\eqref{eq:fixed-source-row-optimum}.

In the finite Gaussian reference regime this dependence is explicit. Let \(f_{ij}=\mathcal N(m_{ij},S_{ij})\) on \(\R^{d}\) with \(S_{ij}\in\mathbb S_{++}^{d}\). Then \(f_{ij}^{(\tau)}=\mathcal N(m_{ij},\tau S_{ij})\) and
\begin{equation}
\log Z_{ij}(\tau)
=\frac d2\left(1-\frac1\tau\right)\log 2\pi
+\frac d2\log\tau
+\frac12\left(1-\frac1\tau\right)\log\det S_{ij},
\label{eq:tempered-gaussian-normalizer}
\end{equation}
whose only source-dependent term is the last. Substituting into Equation~\eqref{eq:tempered-source-row-optimum} gives the tempered Gaussian row
\begin{equation}
\beta_{ij}^{\star(\tau)}
\propto
\pi_{ij}\exp\left[
-\frac{1}{\tau}\KL(\varrho_i\Vert f_{ij})
-\frac12\left(1-\frac1\tau\right)\log\det S_{ij}
\right].
\label{eq:tempered-gaussian-row}
\end{equation}

Equation~\eqref{eq:tempered-gaussian-row} coincides with a plain tempered softmax of the divergences if and only if \(\tau=1\), or \(\det S_{ij}\) does not depend on \(j\), or the declared categorical prior is \(\pi_{ij}Z_{ij}(\tau)^{-1}\) rather than \(\pi_{ij}\). Otherwise the exact tempered model carries an additional per-source log-determinant logit. For \(\tau>1\) its coefficient is positive, so the exact construction discounts sources whose component occupies more volume, beyond whatever the divergence already records. When the component arises as a linear pushforward \(S_{ij}=\Omega_{ij}\Sigma_j\Omega_{ij}^\top\) of a sender covariance, the term reads \(\log\det\Sigma_j+2\log\lvert\det\Omega_{ij}\rvert\). This logit does not cancel inside the divergence, and it is present at every \(\tau\) including \(\tau=1\); a relative entropy between Gaussians is not invariant under a linear pushforward of one argument alone. What does hold is the linked-endpoint statement: under a change of receiver frame \(g_i\) both arguments of the divergence push forward by \(g_i\), so the divergence itself is unchanged, while the additional logit shifts by the source-free constant \(-(1-\tau^{-1})\log\lvert\det g_i\rvert\), which cancels under row normalization. The logit is therefore a property of the tempering and of the sender geometry, and it is consistent with gauge covariance without being removed by it.

\section{A source-label augmented model and its exact attention row}
\label{sec:source-label-augmented-model}

Chapter~\ref{ch:pifb2-crosswalk} records that a fixed normalized model carrying a declared source label can possess an exact ELBO once mutual compatibility of its conditionals and its global normalization are proved. This section carries out that construction, states the normalization condition it actually requires, and identifies the resulting label coordinate. The construction enlarges the latent inventory and is therefore not an equality of functionals on the original shared agent-state variables; the scope statements of Chapter~\ref{ch:pifb2-crosswalk} continue to apply unchanged.

Work at one design point and suppress the index \(a\), with state coordinates \(y_i\in\R^{d}\) for \(i\in V\). Adjoin to the state inventory a label \(j_i\in\mathcal J_i\subseteq V\) for each agent, and declare the normalized factors
\begin{equation}
p(J)=\prod_{i\in V}\pi_{ij_i},
\qquad
p(y_i\given y_{j_i},j_i)
=\mathcal N\left(y_i;\Omega_{ij_i}y_{j_i},R_{ij_i}\right),
\label{eq:source-label-factors}
\end{equation}
with \(\Omega_{ij}\in GL(d,\R)\) fixed transports and \(R_{ij}\in\mathbb S_{++}^{d}\) fixed link covariances. Both are generative parameters and neither reads a live recognition marginal.

Global normalization is not automatic, and the obstruction is combinatorial rather than analytic.

\paragraph{Proposition (label-graph normalization).}
\label{prop:label-graph-normalization}
For every finite \(V\) and every assignment \(i\mapsto j_i\) giving each agent exactly one parent, the induced parent map has a cycle, so the product of the factors in Equation~\eqref{eq:source-label-factors} is not a directed factorization of a probability density. If the source sets \(\mathcal J_i\) are restricted so that the parent relation admits a topological order, with the agents minimal in that order carrying declared normalized priors \(p_i(y_i)\) in place of a link factor, then for every label configuration \(J\) the product is a directed factorization of a normalized joint and its partition function equals one identically \citep{Lauritzen1996}.

Excluding only the self edge is insufficient, since two-cycles survive that exclusion. A strictly ordered source mask, of which a causal mask is the canonical instance, is sufficient. The failure in the unordered case is not a matter of an unknown constant. Under a cocycle transport of the form \(\Omega_{ij}=U_iU_j^{-1}\), so that \(\Omega_{ji}=\Omega_{ij}^{-1}\), a reciprocal pair of link factors has a singular assembled precision for every \(R\in\mathbb S_{++}^{d}\), because both factors express the same linear agreement relation and the joint is flat along it. That statement and its kernel vector are proved in Appendix~\ref{sec:tempered-and-labeled-source-proofs}. Adjoining a proper self-prior to each agent restores positive definiteness but converts the model into an undirected specification whose partition function depends on \(J\), which is treated at the end of this section.

Under the ordered mask the model of Equation~\eqref{eq:source-label-factors} is a fixed normalized joint on the enlarged inventory, so the state ELBO of Chapter~\ref{ch:structured-recognition-elbo} and the coordinate equation of Equation~\eqref{eq:cavi-normalized-coordinate-density} both apply. Take the product recognition family
\begin{equation}
Q=\prod_{i\in V}q_i(y_i)\prod_{i\in V}\beta_i(j_i),
\qquad
q_i=\mathcal N(\mu_i,\Sigma_i),
\qquad
\sum_{j\in\mathcal J_i}\beta_{ij}=1.
\label{eq:source-label-recognition-family}
\end{equation}

\paragraph{Proposition (exact label coordinate).}
\label{prop:exact-label-coordinate}
For the family of Equation~\eqref{eq:source-label-recognition-family} the exact coordinate maximizer of the state ELBO in the factor \(\beta_i\) is
\begin{equation}
\beta_{ij}^\star\propto\pi_{ij}\exp\left(-E_{ij}\right),
\qquad
E_{ij}:=-\E_{q_i\otimes q_j}\left[\log\mathcal N\left(y_i;\Omega_{ij}y_j,R_{ij}\right)\right],
\label{eq:source-label-row-optimum}
\end{equation}
and, writing \(\Delta_{ij}=\mu_i-\Omega_{ij}\mu_j\) and \(S_{ij}=\Omega_{ij}\Sigma_j\Omega_{ij}^\top\),
\begin{equation}
E_{ij}
=\frac12\left[
d\log 2\pi
+\log\det R_{ij}
+\Delta_{ij}^\top R_{ij}^{-1}\Delta_{ij}
+\operatorname{tr}\left(R_{ij}^{-1}\Sigma_i\right)
+\operatorname{tr}\left(R_{ij}^{-1}S_{ij}\right)
\right].
\label{eq:source-label-edge-energy}
\end{equation}
If in addition the link covariance is tied to the pushforward sender covariance, \(R_{ij}=S_{ij}\), then
\begin{equation}
E_{ij}-\KL\left(q_i\middle\Vert(\Omega_{ij})_{\#}q_j\right)
=\frac d2\log 2\pi+\frac12\log\det\Sigma_i+d,
\label{eq:source-label-tied-offset}
\end{equation}
which carries no source index, so in that case
\begin{equation}
\beta_{ij}^\star
=\frac{\pi_{ij}\exp\left[-\KL\left(q_i\middle\Vert(\Omega_{ij})_{\#}q_j\right)\right]}
{\displaystyle\sum_{\ell\in\mathcal J_i}
\pi_{i\ell}\exp\left[-\KL\left(q_i\middle\Vert(\Omega_{i\ell})_{\#}q_\ell\right)\right]}.
\label{eq:source-label-tied-row}
\end{equation}
The proofs are in Appendix~\ref{sec:tempered-and-labeled-source-proofs}.

Equation~\eqref{eq:source-label-tied-row} is a divergence-scored attention row obtained as an exact variational coordinate of a fixed normalized joint, rather than as the value of an engineered simplex minimization. Its scope should be read precisely. The sender mean \(\mu_j\) is a free recognition quantity throughout and is never replaced by a template, which is a strictly weaker requirement than the frozen-source hypothesis of Equation~\eqref{eq:fixed-source-row-joints}. The sender covariance is not equally free, because the tie \(R_{ij}=S_{ij}\) constrains a generative parameter to agree with \(\Sigma_j\). Within a single coordinate update the generative model is fixed and Equation~\eqref{eq:source-label-tied-row} is exact; if \(R_{ij}\) is refreshed to follow a moved \(\Sigma_j\), the model changes between updates, which is the same caution recorded after Equation~\eqref{eq:fixed-source-row-optimum}. Holding \(R_{ij}\) genuinely fixed instead leaves Equation~\eqref{eq:source-label-tied-row} valid at the operating point where \(\Sigma_j\) agrees with the declared link covariance, and elsewhere the row is Equation~\eqref{eq:source-label-edge-energy} with untied \(R_{ij}\). Untied, the edge energy is a Mahalanobis term in the model's own link metric together with separate dispersion penalties on the receiver and on the pushforward sender. It should not be said that the tied case is a specialization of a family that also contains the deployed rule at a nonunit temperature: Section~\ref{sec:link-covariance-status} shows that no choice of \(R_{ij}\) delivers a divergence-scored row when \(\tau\neq1\).

Two further remarks bound the construction. First, adjoining a distinguished label value with \(p(y_i\given j_i=\varnothing)=p_i(y_i)\) keeps the model in the same class and contributes the slot energy
\begin{equation}
-\E_{q_i}\left[\log p_i\right]=\KL(q_i\Vert p_i)+H(q_i),
\label{eq:source-label-self-slot}
\end{equation}
so a self term and the source terms become competing slots of one simplex, and moving a self term inside that simplex costs exactly one receiver entropy. Second, if the ordered mask is dropped in favor of a proper self-prior anchor, the resulting specification is undirected. Its partition function then depends on the label configuration, so \(-\log Z(J)\) enters the label objective and does not separate across receivers; in that route Equation~\eqref{eq:source-label-tied-row} holds only under the further declared approximation that this coupling is neglected.

\section{The status of the link covariance}
\label{sec:link-covariance-status}

Equation~\eqref{eq:source-label-tied-row} rests on the tie \(R_{ij}=S_{ij}\), and that hypothesis can be graded rather than left standing as a convenience. Three results fix its status, and together they place it as a declared modeling postulate: a constrained-covariance choice of the kind that model-based clustering compares between candidate models rather than justifies as an estimator.

The first result is that the tie and the unit temperature are one hypothesis rather than two. Require the offset \(E_{ij}-\tau^{-1}\KL(q_i\Vert\Omega_{ij\#}q_j)\) to be free of the source index as an identity in the sender mean, which is unconstrained and which sweeps \(\R^d\) because \(\Omega_{ij}\) is invertible. Only the quadratic form in \(\Delta_{ij}\) carries that dependence, and its coefficient \(\tfrac12(R_{ij}^{-1}-\tau^{-1}S_{ij}^{-1})\) vanishes identically only when
\begin{equation}
R_{ij}=\tau S_{ij}.
\label{eq:link-covariance-forced-scale}
\end{equation}
Substituting Equation~\eqref{eq:link-covariance-forced-scale} back, the surviving term varies with \(\log\det S_{ij}\) at the rate \((\tau-1)/(2\tau)\), which vanishes only at \(\tau=1\). Fixing the scale therefore fixes the temperature. A corollary worth stating separately is that \(R_{ij}=cS_{ij}\) reproduces Equation~\eqref{eq:tempered-gaussian-row} exactly at \(\tau=c\), log-determinant logit included, so the link-covariance scale of the generative model and the softmax temperature of the recognition row are the same parameter seen from two sides. The consequence for the crosswalk of Chapter~\ref{ch:pifb2-crosswalk} is sharp: a divergence-scored row at a nonunit temperature is not the exact label coordinate of any model in this family, whether the link covariance is tied, scaled, or left free.

The second result is that the tie is not a stationary point of the same objective, so it cannot be discharged by appeal to the slow timescale.

\paragraph{Proposition (M-step optimum of the link covariance).} Since \(\Delta_{ij}^\top R^{-1}\Delta_{ij}=\operatorname{tr}(R^{-1}\Delta_{ij}\Delta_{ij}^\top)\), the edge energy of Equation~\eqref{eq:source-label-edge-energy} depends on \(R\) only through
\begin{equation}
E_{ij}(R)=\tfrac12\left[d\log2\pi+\log\det R+\operatorname{tr}\left(R^{-1}C_{ij}\right)\right],
\qquad
C_{ij}:=\Sigma_i+S_{ij}+\Delta_{ij}\Delta_{ij}^\top.
\label{eq:link-covariance-sufficient-statistic}
\end{equation}
Maximizing the evidence lower bound over \(R\) at fixed recognition factors has the unique solution \(R^\star_{ij}=C_{ij}\). Hence
\begin{equation}
R^\star_{ij}-S_{ij}=\Sigma_i+\Delta_{ij}\Delta_{ij}^\top\succ0
\label{eq:link-covariance-mstep-gap}
\end{equation}
for every configuration of the recognition family, so \(R_{ij}=S_{ij}\) is never stationary, and the evidence lower bound the tie forgoes on each edge is exactly the Stein loss
\begin{equation}
E_{ij}(S_{ij})-E_{ij}(R^\star_{ij})=\tfrac12\left[\operatorname{tr}\left(S_{ij}^{-1}C_{ij}\right)-\log\det\left(S_{ij}^{-1}C_{ij}\right)-d\right]>0.
\label{eq:link-covariance-stein-deficit}
\end{equation}

Equality in Equation~\eqref{eq:link-covariance-mstep-gap} would require \(\Sigma_i=0\) together with \(\Delta_{ij}=0\), a point mass at exact gauge consensus, which lies outside the recognition family. The two terms the tie omits have plain readings. The receiver covariance \(\Sigma_i\) is the mean-field residual spread, present because the quantity being predicted is itself a belief rather than an observation; the outer product \(\Delta_{ij}\Delta_{ij}^\top\) is the squared prediction error that an M-step exists to absorb. Equation~\eqref{eq:link-covariance-stein-deficit} therefore carries an interpretation as well as a number: the divergence-scored row prices its sources as though the receiver were already certain and already in agreement.

The third result closes the remaining escape route, and it matters because relaxing the mean-field pair factorization is an independently motivated direction rather than an artificial one.

\paragraph{Proposition (stationarity and divergence scoring are exclusive).} Let the pair recognition factor be jointly Gaussian with cross-covariance \(C:=\operatorname{Cov}(y_i,y_j)\), so that the residual second moment becomes \(W_{ij}=\Sigma_i+S_{ij}-C\Omega_{ij}^\top-\Omega_{ij}C^\top+\Delta_{ij}\Delta_{ij}^\top\). Then the tie is an M-step fixed point if and only if
\begin{equation}
C\Omega_{ij}^\top+\Omega_{ij}C^\top=\Sigma_i+\Delta_{ij}\Delta_{ij}^\top,
\label{eq:link-covariance-correlated-condition}
\end{equation}
which is solvable and admissible as a joint law on a nonempty set, while the exact label coordinate is a divergence-scored row if and only if \(\operatorname{tr}(S_{ij}^{-1}(C\Omega_{ij}^\top+\Omega_{ij}C^\top))\) does not depend on the source. Because Equation~\eqref{eq:link-covariance-correlated-condition} determines that symmetric combination uniquely, the two requirements hold together only on a nongeneric set. The mean-field factor \(C=0\) is the unique member of the family that yields the row, and it is exactly the member at which the tie fails to be stationary.

A correlated recognition family therefore buys stationarity or divergence scoring, never both. Taken with the first result, this bounds what the construction of Section~\ref{sec:source-label-augmented-model} establishes. The divergence-scored row at unit temperature is realizable as the exact coordinate of a declared normalized joint under an ordered source mask, and \(R_{ij}=S_{ij}\) is the unique link covariance realizing it; but that uniqueness is a statement about the target rule, since the model was selected by requiring it, and it is not a derivation of the rule from independent principles. The proofs are in Appendix~\ref{sec:link-covariance-status-proofs}.

One further remark forestalls a natural reading of the tie. Setting the link covariance to the transported sender covariance does not make the link reproduce the transported sender. Marginalizing the sender's belief through the link gives predictive covariance \(R_{ij}+S_{ij}=2S_{ij}\) under the tie, so the channel is twice as dispersed as the belief it is meant to carry, and the choice that would reproduce the transported belief exactly is \(R_{ij}=0\), outside the positive-definite cone.

\paragraph{Scope of the Gaussian assumption.} The three results above are often read as artifacts of assuming Gaussian beliefs, and they are not. The governing statement is that the M-step reads the recognition family only through the \emph{link's} expected sufficient statistic. For the Gaussian link of Equation~\eqref{eq:source-label-factors}, \(-\log\mathcal N(y_i;\Omega_{ij}y_j,R)\) is affine in \(zz^\top\) for \(z:=y_i-\Omega_{ij}y_j\), so that statistic is the residual second moment and no feature of the recognition factors above second order can enter. Consequently Equations~\eqref{eq:link-covariance-mstep-gap} and \eqref{eq:link-covariance-stein-deficit} hold for any product recognition family whose factors have a mean and a finite nonsingular residual second moment, whether or not those factors are Gaussian, and the gap is positive definite whenever \(\Sigma_i\) is. What the argument does not remove is the product hypothesis, and the correct weakening is not independence but \(\operatorname{sym}(C\Omega_{ij}^\top)=0\) for \(C:=\operatorname{Cov}(y_i,y_j)\), a condition strictly weaker than \(C=0\) once \(d\geq2\). Dependence, not higher-order moment structure, is the only recognition-side freedom that moves the optimum, which is why the exclusivity of the second proposition rather than this observation is what closes the structured route. A factor with no finite second moment is not a hard case but an excluded one: \(z^\top R^{-1}z\geq\lambda_{\min}(R^{-1})\lVert z\rVert^2\) makes the edge energy infinite on all of \(\mathbb S_{++}^d\) simultaneously, so such a factor lies outside the domain of the functional and the Gaussian link imposes the moment restriction by itself.

Relaxing Gaussianity of the beliefs in fact costs the divergence-scored row rather than recovering the tie. Under the tie the offset is \(H(q_i)+d/2+D_{ij}\) with \(D_{ij}=\E_{q_i}[\log(\Omega_{ij\#}q_j/\Pi_{\mathcal G}\Omega_{ij\#}q_j)]\), the discrepancy between the transported sender and its own Gaussian moment match, so Equation~\eqref{eq:source-label-tied-row} survives exactly when the receiver cannot separate each source from that match. Among elliptical families with a common generator closed under linear pushforward, and with the source geometry free, the offset admits a source-free gauge-covariant link-covariance rule if and only if the family is Gaussian, in which case that rule is uniquely the tie; for minimal exponential families the same argument leaves only sufficient statistics inside \(\operatorname{span}\{y,yy^\top\}\). The Gaussian recognition family is therefore not a simplification that happens to work here but the one setting in which this construction closes.

The assumption that does carry weight is Gaussianity of the \emph{link}. For an elliptical link \(\lvert R\rvert^{-1/2}g(z^\top R^{-1}z)\) the stationarity condition becomes the weighted fixed point \(R^\star=\E[w(Q_{R^\star})zz^\top]\) with \(w=-2(\log g)'\), and contracting the tie condition against \(S_{ij}^{-1}\) gives the necessary scalar requirement \(\E[\psi(Q_{S_{ij}})]=d\) with \(\psi(u)=uw(u)\). The Gaussian case is exactly \(\psi=\mathrm{id}\), which returns \(d+\operatorname{tr}(S_{ij}^{-1}(\Sigma_i+\Delta_{ij}\Delta_{ij}^\top))=d\) and forces the degenerate configuration; for a heavy-tailed generator \(\psi\) is bounded and the requirement acquires solutions, so the tie is attainable on a nonempty set of edge configurations. It does not become an identity, since a nonzero mean residual obstructs it and demanding it identically forces \(g(u)\propto u^{-d/2}\), which is not normalizable, and where it does hold the exact row is a robust score rather than a Gaussian divergence. The exclusivity therefore migrates to the non-Gaussian link rather than dissolving there, and that case is left open.
